%% ma: L12-B03-0009
%% lop: 12
%% chuong: 1
%% bai: 3
%% dang: 12.3.7
%% loai: TN4
%% muc: TH
%% dap_an: "C"
%% nguon: "Ảnh giáo viên gửi 11/10/2026 (Đề kiểm tra giữa HK1 Toán 12, THPT Việt Đức 2026-2027, mã đề 0103), Phần I Câu 9"
%% kien_thuc: "Bài 3 (tiệm cận đứng, tiệm cận ngang của đồ thị hàm số)"
%% trang_thai: cho-duyet
%% ly_do: "Dạng mới đề xuất, chờ giáo viên duyệt."
\begin{ex}
Đồ thị hàm số nào sau đây có đúng $2$ đường tiệm cận?
\choice{$y=\dfrac{x^2+x}{2}$}{$y=x^3-2x^2-1$}{\True $y=\dfrac{1-x}{x+2}$}{$y=x+1$}
\loigiai{Các hàm số đa thức $y=\dfrac{x^2+x}2$, $y=x^3-2x^2-1$, $y=x+1$ không có tiệm cận đứng và tiệm cận ngang (đồ thị của chúng cũng không có tiệm cận xiên). Hàm số $y=\dfrac{1-x}{x+2}$ có tiệm cận đứng $x=-2$ và tiệm cận ngang $y=-1$, tức đúng $2$ đường tiệm cận.}
\end{ex}
